% Generated draft: requires mathematical and source review.
\documentclass[11pt]{article}
\usepackage[margin=25mm]{geometry}
\usepackage{fontspec}
\setmainfont{Latin Modern Roman}
\usepackage{amsmath,amssymb,mathtools,unicode-math}
\setmathfont{Latin Modern Math}
\usepackage{enumitem,needspace,xurl,hyperref}
\setlength{\parindent}{0pt}
\setlength{\parskip}{5pt}
\setlength{\emergencystretch}{3em}
\newcommand{\N}{\mathbb{N}}
\newcommand{\Z}{\mathbb{Z}}
\newcommand{\Q}{\mathbb{Q}}
\newcommand{\R}{\mathbb{R}}
\newcommand{\C}{\mathbb{C}}
\newcommand{\F}{\mathbb{F}}
\newcommand{\A}{\mathbb{A}}
\newcommand{\PP}{\mathbb{P}}
\newcommand{\E}{\mathbb{E}}
\newcommand{\eps}{\varepsilon}
\newcommand{\dd}{\,\mathrm{d}}
\begin{document}
\section*{UnsolvedMath ID: 30000674}
\textbf{Source problem code:} catalogue entry
\subsection{Definitions and mathematical statement}
Let \(G\) be a finitely generated group, let \(\alpha\in\operatorname{Aut}(G)\), and fix \(g\in G\). Its forward automorphism orbit is \((\alpha^n(g))_{n\geq0}\). Define the period, when it exists, by
\[
\operatorname{per}_\alpha(g)=\min\{m\geq1:\alpha^m(g)=g\}.
\]
The first task is to describe possible periods or bounds for periodic elements. The source does not say whether a bound may depend on \(g\), on \(\alpha\), on \(G\), or on a numerical group parameter, so no uniform-period theorem is asserted.

For the nonperiodic case, fix a finite symmetric generating set \(S\) and define
\[
L_{\alpha,g,S}(n)=|\alpha^n(g)|_S.
\]
The second task is to describe its growth. Changing finite generating sets changes word lengths only up to multiplicative comparison, but the source does not identify the desired growth equivalence, exact rate, or class of groups. These are lengths of actual elements; length minimized within conjugacy classes would be a different quantity and is not substituted.

The third task asks whether the sequence converges in a suitable sense, perhaps after replacing \(\alpha\) by a positive power. This would concern \((\alpha^{mn}(g))_n\) for some \(m\geq1\), in a specified topology, boundary, or compactification. With the discrete topology on \(G\), convergence means eventual constancy and gives no nontrivial convergence notion for a nonperiodic orbit. An arbitrary finitely generated group has no single boundary supplied by the statement, so a hyperbolic boundary, profinite completion, or other target cannot be chosen silently. The dependence of \(m\) on \(G,\alpha,g\), and any uniformity over all elements, is likewise unspecified.

The source asks a broad three-part asymptotic programme for finitely generated groups, not an assertion that all orbits have one prescribed period, growth rate, or limit. No restriction to free, hyperbolic, or other special groups is added. The supplied literature triage checked 2026-08-21 reports no authoritative complete resolution of the extracted collection, without individual theorems for its parts.
\subsection{Short English statement}
Describe periods, word-length growth, and meaningful limits of automorphism orbits in finitely generated groups, allowing an iterate of the automorphism. The required uniformity, growth comparison, and convergence space are missing from the broad question.
\subsection{Sources}
\begin{itemize}
\item Local UnsolvedMath registry, permanent ID 30000674, catalogue entry, \emph{Asymptotic Dynamics of Group Automorphism Orbits}.
\item \emph{Model Theory and Groups} (2007), Oberwolfach Reports, pp. 83--138, DOI \url{https://doi.org/10.4171/owr/2007/02}.
\item Supplied source-verification claim and literature triage checked 2026-08-21 cite only the report; no separate period, growth, or boundary-convergence theorem is supplied.
\end{itemize}
\end{document}
